\section{Trace and products}\label{sec:trace-products}

The trace is centered in the phase already selected in Section~\ref{sec:phases}. Its short-distance logarithm will distinguish the resulting theory from finite Wick comparisons.

\begin{definition}[Centered trace]\label{def:D16}
For the tensor of Theorem~\ref{prop:ST}, put
\begin{align*}
\tau&=\sum_{a=1}^4T_{aa},&e_{\rm tr}&=S^{\xi_*}(\tau),&\Theta&=\tau-e_{\rm tr}\unitfield{4},\\
\tracetwopoint(r)&=S_2^{\xi_*}(\Theta,\Theta;-re_4/2,re_4/2),&&\traceamplitude&=\frac{3\dimG}{\pi^4}\,.
\end{align*}
The one-point value is position independent. Centering removes the vacuum contribution in the selected phase.
\end{definition}

The normalized stress has a nonzero centered trace. The next limit determines both its leading power and its logarithmic factor.

\begin{theorem}[Trace law]\label{prop:TR}
The selected field $\Theta$ is nonzero, real and $O(4)$ invariant; its relativistic field is a Hermitian Lorentz scalar. For every $f\in\mathcal S$,
\begin{equation*}
\langle\Omega,\Phi(\iota_4\Theta,f)\Omega\rangle=0,\qquad \lim_{r\downarrow0}r^8\log^2(1/r)\tracetwopoint(r)=\traceamplitude>0\,.
\end{equation*}
\end{theorem}

The proof is in Section~\ref{sec:partII-I}, I.1.1--I.1.3. The degree-four comparison combines the scalar coefficient $d_2=-b_0/2$ with the running-coupling slope $\slope=2b_0$. Its leading trace coefficient is therefore $4d_2/(\slope\log(1/r))=-1/\log(1/r)$. Squaring it gives the logarithmic factor; the free curvature-square correlation gives $\traceamplitude$. The comparison uses adjustable stopping levels and fourth-order information beyond Proposition~\ref{prop:MC}. It proves the limit along all positive radii.

The same proof identifies centering on the word domain:
\begin{equation*}
\Phi(\iota_4\Theta,f)=\Phi(\iota_4\tau,f)-e_{\rm tr}\left(\int f\right)\mathrm{Id}_{\mathcal H}\,.
\end{equation*}
It also constructs the unique positive tempered-growth spectral measure $\tracespectrum$, supported in $\forwardcone$, for which
\begin{equation*}
d\mu_{\Phi(\iota_4\Theta,f)\Omega}=|\fourierplus{f}|^2\,d\tracespectrum,\qquad \tracetwopoint(r)=\int e^{-r\mathbf p^4}\,d\tracespectrum(\mathbf p)\quad(r>0)\,.
\end{equation*}
Here $\fourierplus{f}(\mathbf p)=\int e^{i\mathbf p\cdot x}f(x)\,dx$. The measure's Lorentz invariance and its high-energy limits against compactly supported continuous tests and Schwartz tests are stated in I.1.2--I.1.3.

\begin{definition}[Finite Wick profiles]\label{def:D17}
For $r>0$, $m_{\rm sh}\ge0$, integers $j_h,k_E\ge0$ and continuous real $h_{\rm mom}$ on $\R^3$ with $h_{\rm mom}(\lambda\mathbf k)=\lambda^{j_h}h_{\rm mom}(\mathbf k)$ for $\lambda>0$, set
\begin{align*}
\omega_{m_{\rm sh}}(\mathbf k)&=\sqrt{m_{\rm sh}^2+|\mathbf k|^2},\\
K_{m_{\rm sh},k_E,h_{\rm mom}}(r)&=\int_{\R^3}\frac{\omega_{m_{\rm sh}}(\mathbf k)^{k_E}h_{\rm mom}(\mathbf k)}{2\omega_{m_{\rm sh}}(\mathbf k)}e^{-r\omega_{m_{\rm sh}}(\mathbf k)}\,d^3\mathbf k\,.
\end{align*}
A finite Wick profile is a finite complex linear combination of finite products of these kernels. A zero denominator at zero momentum gives integrand zero. This is the Euclidean comparison class.
\end{definition}

Finite products and sums of these kernels cannot reproduce the inverse logarithm in Theorem~\ref{prop:TR}.

\begin{proposition}[Finite-profile exclusion]\label{prop:NF}
For every finite Wick profile $W_{\rm Wick}$ and every $\delta>0$, there is $0<r<\delta$ such that
\begin{equation*}
\tracetwopoint(r)\ne W_{\rm Wick}(r)\,.
\end{equation*}
\end{proposition}

The proof is in Section~\ref{sec:partII-I}, I.2.1. The mass-shell kernels have expansions to every finite order in integer powers of $r$ and nonnegative powers of $\log r$. Finite products and sums therefore give, to bounded error, a finite sum of monomials $r^a(\log r)^k$ with $a\in\Z$ and $k\in\N$. Multiplication by $r^8\log^2(1/r)$ makes the terms with $a>-8$ vanish. The remaining nonzero leading term diverges, so any finite limit is zero. The positive limit in Theorem~\ref{prop:TR} therefore excludes equality on every sufficiently small interval. I.2.2--I.2.4 give the polynomial mass-shell and Fock countertests, with their comparison classes defined there.

The product expansion is stated for the Euclidean tables on their common quotient fields. Write $\zeta\in\phases\sqcup\{\mathrm{mix}\}$, with $S^{\rm mix}=S$, and retain degree tags by $\mathbb F=\bigsqcup_dF_d$, $\deg(d,A)=d$. Choose finite channel sets $\channels{j}$ and real fields $O_{j,\kappa}\in F_j$ forming complex bases of conjugation-invariant complements of $F_{j-1}$ in $F_j$, with $F_{-1}=0$. Section~\ref{sec:partII-J}, J.1.1 constructs these fields and coefficient functions $c_{j,\kappa}$ once, before spectators, states and truncations. The degree-zero vector spans the unit line; it need not equal the unit. For finite cluster and spectator index sets $I_{\rm cl},I_{\rm sp}$, the positions are $x\in(\R^4)^{I_{\rm cl}}$, $y\in(\R^4)^{I_{\rm sp}}$, with expansion center $o\in\R^4$.

\begin{definition}[Product remainder]\label{def:D18}
For finite index sets $I_{\rm cl},I_{\rm sp}$ and $D\in\N$, cluster and spectator lists $\boldsymbol A\in\mathbb F^{I_{\rm cl}}$, $\boldsymbol B\in\mathbb F^{I_{\rm sp}}$, put $|\boldsymbol A|_{\deg}=\sum_i\deg A_i$ and
\begin{equation*}
\mathcal R_D^\zeta=S^\zeta(\boldsymbol A,x;\boldsymbol B,y)-\sum_{j=0}^D\sum_{\kappa\in\channels{j}}c_{j,\kappa}(\boldsymbol A;x,o)S^\zeta(O_{j,\kappa}(o),\boldsymbol B(y))\,.
\end{equation*}
Use a fixed enumeration of $I_{\rm cl}\sqcup I_{\rm sp}$ for the concatenated correlations. An $(r,M_{\rm geo},\eta)$-cluster satisfies $|x_i-o|\le M_{\rm geo}r$ and $|x_i-x_{i'}|\ge\eta r$ for distinct slots. Spectators are $\delta_{\rm sp}$-clear if $|o-y_j|\ge\delta_{\rm sp}$ and $|y_j-y_{j'}|\ge\delta_{\rm sp}$ for distinct slots. The remainder measures the error through degree $D$.
\end{definition}

One field expansion approximates a cluster in every phase and in the mixture. Its coefficients are independent of all spectators.

\begin{theorem}[Operator product expansion]\label{prop:OP}
For every finite list $\boldsymbol A$, $M_{\rm geo}\ge0$, $\eta,\varepsilon>0$, $j\in\N$ and $\kappa\in\channels{j}$, some $r_0>0$ gives, for $0<r<r_0$ and every $(r,M_{\rm geo},\eta)$-cluster,
\begin{equation*}
|c_{j,\kappa}(\boldsymbol A;x,o)|\le r^{j-|\boldsymbol A|_{\deg}-\varepsilon}\,.
\end{equation*}
For every pair of finite lists $\boldsymbol A,\boldsymbol B$, $M_{\rm geo}\ge0$, $\eta,\delta_{\rm sp},\varepsilon>0$ and $D\in\N$, some $r_1>0$ gives, for $0<r<r_1$, such clusters, $\delta_{\rm sp}$-clear spectators and all states,
\begin{equation*}
|\mathcal R_D^\zeta|\le r^{D+1-|\boldsymbol A|_{\deg}-\varepsilon}\,.
\end{equation*}
\end{theorem}

The proof is in Section~\ref{sec:partII-J}, J.1.1--J.1.2. Finite correlation functionals first extract a field, rotation and conjugation averages impose the symmetries, and recursive subtraction removes lower degrees. The lattice product estimate controls the discarded part. Its compressed-algebra norm bound holds before character evaluation, which gives the same radius threshold for every state. The two thresholds have separate dependencies, and each displayed bound has coefficient one.

J.1.3--J.1.6 give integrated and nested estimates; the nested error is a sum of gains from proper clusters.

Separated product estimates leave the collision diagonals untreated. Contact subtractions use the same underlying prescription, with their own projections, separating correlation functionals and clearance.

\begin{theorem}[Contact extension]\label{prop:CT}
For every state $\zeta$, arity $n$ and list $A_i\in F_{d_i}$, the fixed contact construction gives a distribution $\mathfrak D_{\boldsymbol A}^\zeta$ on $\R^{4n}$ such that
\begin{equation*}
\mathfrak D_{\boldsymbol A}^\zeta(f)=\int f(x)S_n^\zeta(\boldsymbol A;x)\,dx\qquad(f\in C_c^\infty(\Conf{n}))\,.
\end{equation*}
\end{theorem}

The proof is in Section~\ref{sec:partII-J}, J.2.1--J.2.3. Taylor subtraction in the directions normal to collision diagonals is organized by forests of disjoint or nested subsets. The contact projections retain compatible field components, while the construction's lattice bounds make the shell sums converge. That section specifies the real stratum bases, subtraction scales and local finite-order bounds. These contact estimates are additional to the separated OPE bounds.
